% Template for post or presentation abstracts submitted to ILAS 2022
% 1. Fill in the presenter's information (name, affiliation, email
% address) and talk information (title of presentation, abstract).
% 2. Compile to make sure there are no errors.
% 3. Save the .tex file for your abstract with a distinctive name,
% ideally "FamilynameGivenname.tex".
% 4. Upload your .tex file to https://clr.ie/129557
%
% * Please keep your LaTeX commands simple, and avoid the use of
%    user-defined macros.
% * Include details of co-authors and funding acknowledgements after the abstract.
% * Upload only the LaTeX file (with .tex extension).
% * Questions: Contact Niall Madden (Niall.Madden@NUIGalway.ie)

\documentclass[a4paper]{amsart}
\usepackage{amssymb}
\usepackage{hyperref}
\pagestyle{empty}
\date{}

%% IGNORE THE NEXT 4 LINES
\makeatletter %
\def\labelenumi{\theenumi.}
\def\theenumi{\@arabic\c@enumi}
\makeatother
%% STOP IGNORING NOW

\begin{document}

\title{Seventieth anniversary of the conjugate gradient method and what do old papers reveal about our presence}
\author{Zden\v{e}k Strako\v{s}}
\address{ Charles University, Prague}
\email{strakos@karlin.mff.cuni.cz}

\maketitle

In his lecture {\em Why Mathematics?} delivered at the Annual Meeting of the Irish Mathematics Association on October 31, 1966, Cornelius Lanczos said: {\em "The naive optimist who believes in progress and is convinced that today is better than yesterday and in ten years time the world will be infinitely better off than it is today, will come to the conclusion that mathematics (and more generally all the exact sciences) started only about twenty years ago, while the predecessors must have walked around in a kind of limbo of half-digested and improperly conceived ideas."}

This year marks the seventieth anniversary of the paper by Hestenes and Stiefel, which comprehensively described the conjugate gradient method (CG) considered among the most important algorithmic developments of the $20$th century. This paper should be studied together with the three closely related papers by Lanczos published within the period 1950-53. It is worth to notice, e.g., also the papers by Karush and Hayes, published in 1952 and 1954, respectively, as well as a couple of other works of several other authors from the same period.

This contribution will examine how the knowledge present in these seminal papers is reflected in the contemporary literature on CG, and what does it show on the status of the current understanding of the deeply rooted mathematical ideas so beautifully presented many decades ago.

\end{document}


